%=========================================================
% C++ Chapter 8 Lab Handout (Verbatim Korean-safe)
% Topic: Inheritance (상속)
% Compile with XeLaTeX
%=========================================================
\documentclass[11pt,a4paper]{article}

\usepackage[margin=1in]{geometry}
\usepackage{fontspec}
\usepackage{kotex}
\usepackage{hyperref}
\usepackage{enumitem}
\usepackage{fancyhdr}
\usepackage{titlesec}
\usepackage{xcolor}
\usepackage{tcolorbox}
\usepackage{fvextra}

% ---------- Fonts ----------
% If D2Coding isn't available, try: NanumGothicCoding, Noto Sans Mono CJK KR
\setmonofont{D2Coding}[Scale=MatchLowercase]

% ---------- Colors ----------
\definecolor{CodeBg}{RGB}{248,248,248}
\definecolor{CodeFrame}{RGB}{220,220,220}
\definecolor{Accent}{RGB}{40,80,160}

% ---------- Header ----------
\pagestyle{fancy}
\fancyhf{}
\lhead{\textbf{C++ 8장 실습}}
\rhead{\textbf{상속}}
\cfoot{\thepage}

\titleformat{\section}{\Large\bfseries\color{Accent}}{\thesection}{0.6em}{}
\titleformat{\subsection}{\large\bfseries}{\thesubsection}{0.6em}{}

% ---------- Boxes ----------
\newtcolorbox{GoalBox}{
  colback=CodeBg,
  colframe=CodeFrame,
  arc=3pt,
  boxrule=0.8pt,
  left=10pt,right=10pt,top=8pt,bottom=8pt
}
\newcommand{\filetag}[1]{\texttt{\color{Accent}#1}}

% ---------- Code block environment (Unicode-safe) ----------
\DefineVerbatimEnvironment{cppcode}{Verbatim}{
  fontsize=\small,
  frame=single,
  framerule=0.6pt,
  rulecolor=\color{CodeFrame},
  numbers=left,
  numbersep=8pt,
  baselinestretch=1,
  breaklines=true,
  breaksymbolleft=,
  breaksymbolright=,
  commandchars=\\\{\},
  xleftmargin=1.5em
}

\begin{document}

\begin{center}
  {\LARGE \textbf{C++ 8장 실습: 상속}}\\
  \vspace{6pt}
  {\large (inheritance / upcasting \& downcasting / protected / ctor-dtor order / inheritance modes / multiple \& virtual inheritance)}
\end{center}

\vspace{8pt}

\section*{학습 목표}
\begin{GoalBox}
\begin{enumerate}[leftmargin=18pt]
  \item 객체 지향 상속 개념 이해
  \item 상속 선언 및 파생 클래스 객체 구조 이해
  \item 업 캐스팅/다운 캐스팅과 객체 포인터 관계 이해
  \item \texttt{protected} 접근 지정 이해
  \item 상속 관계 생성/소멸 과정 이해
  \item \texttt{public/protected/private} 상속 차이 이해
  \item 다중 상속 선언/활용
  \item 다중 상속 문제점(다이아몬드)과 가상 상속 개념 이해
\end{enumerate}
\end{GoalBox}

\section*{실습 운영 규칙}
\begin{itemize}[leftmargin=18pt]
  \item 각 실습은 \texttt{main()} 포함, \textbf{독립 실행} 형태입니다.
  \item ``(에러 확인)'' 표시는 주석을 풀면 일부러 컴파일 에러가 나도록 만든 학습용 코드입니다.
  \item 다운캐스팅은 위험합니다. 본 장에서는 개념 이해를 위해 \textbf{명시적 캐스팅}을 보여주되, 실제로는 안전한 방식(\texttt{dynamic\_cast})은 \textbf{가상 함수(9장)} 이후에 다루는 것이 좋습니다.
\end{itemize}

\tableofcontents
\newpage

\section{상속 선언과 파생 클래스 기본 동작(예제 8-1: Point $\rightarrow$ ColorPoint)}
\begin{GoalBox}
\textbf{핵심}
\begin{itemize}[leftmargin=18pt]
  \item \texttt{class ColorPoint : public Point} 처럼 선언.
  \item 파생 클래스 객체는 \textbf{기본 클래스 영역 + 파생 클래스 영역}을 함께 가진다.
  \item 파생 클래스는 기본 클래스의 public 멤버 함수를 그대로 사용 가능(\texttt{cp.set(3,4)}).
\end{itemize}
\end{GoalBox}

\noindent\filetag{L1\_PointColorPoint.cpp}
\begin{cppcode}
#include <iostream>
#include <string>
using namespace std;

class Point {
    int x, y; // private: 파생 클래스에서도 직접 접근 불가
public:
    void set(int x, int y) { this->x = x; this->y = y; }
    void showPoint() const { cout << "(" << x << "," << y << ")" << endl; }
};

class ColorPoint : public Point { // Point를 상속
    string color;
public:
    void setColor(string color) { this->color = color; }
    void showColorPoint() const {
        cout << color << ":";
        showPoint(); // Point의 public 멤버 호출
    }
};

int main() {
    ColorPoint cp;
    cp.set(3,4);
    cp.setColor("Red");
    cp.showColorPoint();
}
\end{cppcode}

\section{업 캐스팅: 파생 클래스 포인터 $\rightarrow$ 기본 클래스 포인터}
\begin{GoalBox}
\textbf{핵심}
\begin{itemize}[leftmargin=18pt]
  \item 업 캐스팅: \texttt{Point* pBase = pDer;} 처럼 \textbf{자동} 변환 가능.
  \item 업캐스팅된 포인터로는 기본 클래스의 \textbf{public 멤버만} 접근 가능.
\end{itemize}
\end{GoalBox}

\noindent\filetag{L2\_Upcasting.cpp}
\begin{cppcode}
#include <iostream>
#include <string>
using namespace std;

class Point {
    int x, y;
public:
    void set(int x, int y) { this->x=x; this->y=y; }
    void showPoint() const { cout << "(" << x << "," << y << ")" << endl; }
};

class ColorPoint : public Point {
    string color;
public:
    void setColor(string c) { color = c; }
    void showColorPoint() const { cout << color << ":"; showPoint(); }
};

int main() {
    ColorPoint cp;
    ColorPoint* pDer = &cp;
    Point* pBase = pDer; // 업캐스팅(자동)

    pDer->set(3,4);
    pBase->showPoint(); // OK
    pDer->setColor("Red");
    pDer->showColorPoint();

    // pBase->showColorPoint(); // (에러 확인) 기본 클래스에 없는 멤버
}
\end{cppcode}

\section{다운 캐스팅: 기본 클래스 포인터 $\rightarrow$ 파생 클래스 포인터(주의)}
\begin{GoalBox}
\textbf{핵심}
\begin{itemize}[leftmargin=18pt]
  \item 다운 캐스팅은 \textbf{명시적 캐스팅}이 필요.
  \item 실제로 가리키는 객체가 파생 클래스일 때만 안전.
\end{itemize}
\end{GoalBox}

\noindent\filetag{L3\_Downcasting.cpp}
\begin{cppcode}
#include <iostream>
#include <string>
using namespace std;

class Point {
    int x, y;
public:
    void set(int x, int y) { this->x=x; this->y=y; }
    void showPoint() const { cout << "(" << x << "," << y << ")" << endl; }
};

class ColorPoint : public Point {
    string color;
public:
    void setColor(string c) { color = c; }
    void showColorPoint() const { cout << color << ":"; showPoint(); }
};

int main() {
    ColorPoint cp;
    Point* pBase = &cp; // 업캐스팅

    pBase->set(3,4);
    pBase->showPoint();

    ColorPoint* pDer = (ColorPoint*)pBase; // 다운캐스팅(강제)
    pDer->setColor("Red");
    pDer->showColorPoint();
}
\end{cppcode}

\section{\texttt{protected}: 파생 클래스만 접근 허용}
\begin{GoalBox}
\textbf{핵심}
\begin{itemize}[leftmargin=18pt]
  \item \texttt{protected}는 파생 클래스 내부에서 접근 가능, 외부에서는 접근 불가.
\end{itemize}
\end{GoalBox}

\noindent\filetag{L4\_ProtectedAccess.cpp}
\begin{cppcode}
#include <iostream>
#include <string>
using namespace std;

class Point {
protected:
    int x, y;
public:
    void set(int x, int y) { this->x=x; this->y=y; }
    void showPoint() const { cout << "(" << x << "," << y << ")" << endl; }
};

class ColorPoint : public Point {
    string color;
public:
    void setColor(string c) { color = c; }
    bool equals(const ColorPoint& p) const {
        return (x == p.x && y == p.y && color == p.color);
    }
};

int main() {
    Point p;
    p.set(2,3);
    // p.x = 5; // (에러 확인)

    ColorPoint a; a.set(3,4); a.setColor("Red");
    ColorPoint b; b.set(3,4); b.setColor("Red");
    cout << (a.equals(b) ? "true\n" : "false\n");

    // a.x = 10; // (에러 확인)
}
\end{cppcode}

\section{생성자/소멸자 실행 순서}
\begin{GoalBox}
\textbf{핵심}
\begin{itemize}[leftmargin=18pt]
  \item 생성: 기본 $\rightarrow$ 파생
  \item 소멸: 파생 $\rightarrow$ 기본
\end{itemize}
\end{GoalBox}

\noindent\filetag{L5\_CtorDtorOrder.cpp}
\begin{cppcode}
#include <iostream>
using namespace std;

class A {
public:
    A() { cout << "생성자 A\n"; }
    ~A() { cout << "소멸자 A\n"; }
};
class B : public A {
public:
    B() { cout << "생성자 B\n"; }
    ~B() { cout << "소멸자 B\n"; }
};
class C : public B {
public:
    C() { cout << "생성자 C\n"; }
    ~C() { cout << "소멸자 C\n"; }
};

int main() { C c; }
\end{cppcode}

\section{기본 생성자 없으면 컴파일 오류 + 초기화 리스트로 해결}
\noindent\filetag{L6\_NoDefaultBaseCtor\_Error.cpp}
\begin{cppcode}
#include <iostream>
using namespace std;

class A {
public:
    A(int x) { cout << "매개변수생성자 A " << x << "\n"; }
};

class B : public A {
public:
    B() {
        // (에러 확인) A()가 없어서 컴파일 오류
        cout << "생성자 B\n";
    }
};

int main() { B b; }
\end{cppcode}

\noindent\filetag{L6b\_FixWithInitList.cpp}
\begin{cppcode}
#include <iostream>
using namespace std;

class A {
public:
    A(int x) { cout << "매개변수생성자 A " << x << "\n"; }
};

class B : public A {
public:
    B(int x) : A(x+3) {
        cout << "매개변수생성자 B " << x << "\n";
    }
};

int main() { B b(5); }
\end{cppcode}

\section{TV $\rightarrow$ WideTV $\rightarrow$ SmartTV: 생성자 매개변수 전달}
\noindent\filetag{L7\_TVWideSmart.cpp}
\begin{cppcode}
#include <iostream>
#include <string>
using namespace std;

class TV {
    int size;
public:
    TV() : size(20) {}
    TV(int size) : size(size) {}
    int getSize() const { return size; }
};

class WideTV : public TV {
    bool videoIn;
public:
    WideTV(int size, bool videoIn) : TV(size), videoIn(videoIn) {}
    bool getVideoIn() const { return videoIn; }
};

class SmartTV : public WideTV {
    string ipAddr;
public:
    SmartTV(string ip, int size) : WideTV(size, true), ipAddr(ip) {}
    string getIpAddr() const { return ipAddr; }
};

int main() {
    SmartTV htv("192.0.0.1", 32);
    cout << "size=" << htv.getSize() << "\n";
    cout << "videoIn=" << boolalpha << htv.getVideoIn() << "\n";
    cout << "IP=" << htv.getIpAddr() << "\n";
}
\end{cppcode}

\section{상속 접근 지정자: public/protected/private 상속}
\noindent\filetag{L8\_InheritanceModes.cpp}
\begin{cppcode}
#include <iostream>
using namespace std;

class Base {
public:
    void pub() { cout << "Base::pub\n"; }
protected:
    void prot() { cout << "Base::prot\n"; }
};

class D1 : public Base {
public:
    void test() { pub(); prot(); }
};

class D2 : protected Base {
public:
    void test() { pub(); prot(); }
};

class D3 : private Base {
public:
    void test() { pub(); prot(); }
};

int main() {
    D1 a; a.pub();
    // D2 b; b.pub(); // (에러 확인)
    // D3 c; c.pub(); // (에러 확인)
}
\end{cppcode}

\section{다중 상속: Calculator(Adder + Subtractor)}
\noindent\filetag{L9\_MultipleInheritanceCalculator.cpp}
\begin{cppcode}
#include <iostream>
using namespace std;

class Adder {
protected:
    int add(int a, int b) { return a + b; }
};
class Subtractor {
protected:
    int minus(int a, int b) { return a - b; }
};

class Calculator : public Adder, public Subtractor {
public:
    int calc(char op, int a, int b) {
        switch(op) {
        case '+': return add(a,b);
        case '-': return minus(a,b);
        default:  return 0;
        }
    }
};

int main() {
    Calculator c;
    cout << "2 + 4 = " << c.calc('+',2,4) << "\n";
    cout << "100 - 8 = " << c.calc('-',100,8) << "\n";
}
\end{cppcode}

\section{다이아몬드 문제와 가상 상속(개념)}
\noindent\filetag{L10\_DiamondVirtualInheritance.cpp}
\begin{cppcode}
#include <iostream>
using namespace std;

class A {
public:
    int x;
    A() : x(1) {}
};

class B : public A {};
class C : public A {};
class D_bad : public B, public C {};

class Bv : virtual public A {};
class Cv : virtual public A {};
class D_good : public Bv, public Cv {};

int main() {
    D_good d;
    d.x = 10;
    cout << d.x << "\n";

    // D_bad e;
    // e.x = 10; // (에러 확인) 어떤 A::x 인지 모호
}
\end{cppcode}

\end{document}
